\let\R\relax
\newcommand{\R}{\mathbb{R}}
\let\N\relax
\newcommand{\N}{\mathbb{N}}
\let\Z\relax
\newcommand{\Z}{\mathbb{Z}}
\let\Q\relax
\newcommand{\Q}{\mathbb{Q}}
\let\A\relax
\newcommand{\A}{\mathbb{A}}
\let\F\relax
\newcommand{\F}{\mathbb{F}}
\let\G\relax
\newcommand{\G}{\mathbb{G}}
\let\Gal\relax
\newcommand{\Gal}{\operatorname{Gal}}
\let\Spec\relax
\newcommand{\Spec}{\operatorname{Spec}}
\let\Hom\relax
\newcommand{\Hom}{\operatorname{Hom}}
\let\Ext\relax
\newcommand{\Ext}{\operatorname{Ext}}
\let\Tor\relax
\newcommand{\Tor}{\operatorname{Tor}}
\let\Frob\relax
\newcommand{\Frob}{\operatorname{Frob}}
\let\Cor\relax
\newcommand{\Cor}{\operatorname{Cor}}
\let\Inf\relax
\newcommand{\Inf}{\operatorname{Inf}}
\let\Ver\relax
\newcommand{\Ver}{\operatorname{Ver}}
\let\Def\relax
\newcommand{\Def}{\operatorname{Def}}
\let\Sel\relax
\newcommand{\Sel}{\operatorname{Sel}}
\let\Br\relax
\newcommand{\Br}{\operatorname{Br}}
\let\Pic\relax
\newcommand{\Pic}{\operatorname{Pic}}
\let\End\relax
\newcommand{\End}{\operatorname{End}}
\let\Aut\relax
\newcommand{\Aut}{\operatorname{Aut}}
\let\ord\relax
\newcommand{\ord}{\operatorname{ord}}
\let\coker\relax
\newcommand{\coker}{\operatorname{coker}}
\let\Ann\relax
\newcommand{\Ann}{\operatorname{Ann}}
\let\LT\relax
\newcommand{\LT}{\operatorname{LT}}
\let\rig\relax
\newcommand{\rig}{\operatorname{rig}}
\let\an\relax
\newcommand{\an}{\operatorname{an}}
\let\Sp\relax
\newcommand{\Sp}{\operatorname{Sp}}
\let\SU\relax
\newcommand{\SU}{\operatorname{SU}}
\let\NS\relax
\newcommand{\NS}{\operatorname{NS}}
\let\GL\relax
\newcommand{\GL}{\operatorname{GL}}
\let\Ind\relax
\newcommand{\Ind}{\operatorname{Ind}}
\let\ST\relax
\newcommand{\ST}{\operatorname{ST}}
\let\Sym\relax
\newcommand{\Sym}{\operatorname{Sym}}
\let\vol\relax
\newcommand{\vol}{\operatorname{vol}}
\let\inv\relax
\newcommand{\inv}{\operatorname{inv}}
\let\ab\relax
\newcommand{\ab}{\mathrm{ab}}
\let\Sha\relax
\newcommand{\Sha}{\text{\fontencoding{T2A}\selectfont\CYRSH}}

\chapter{John Tate (2002/03)}

\begin{quote}
\it ``John Tate's work has profoundly shaped modern number theory and arithmetic geometry. His creation of Tate cohomology, his thesis on Fourier analysis in number fields, his discovery of the Tate--Shafarevich group, his conjecture on algebraic cycles, and his foundational contributions to rigid analytic geometry, formal groups, and the arithmetic of elliptic curves have opened entire new fields of investigation. His ideas, characterized by their depth, elegance, and lasting influence, permeate nearly every corner of algebraic number theory.'' --- Wolf Prize Citation
\end{quote}

John Torrence Tate Jr.\ (March 13, 1925 -- October 16, 2019) was one of the most influential mathematicians of the 20th century. His work transformed algebraic number theory, arithmetic geometry, and the theory of automorphic forms. He received the Wolf Prize in Mathematics in 2002/03 (shared with Mikio Sato) for his fundamental contributions to number theory and related fields. In 2010, he was awarded the Abel Prize for his immense and lasting impact on mathematics.

Tate's mathematical legacy is vast. It includes the creation of Tate cohomology (the cohomology theory of finite groups), the adelic approach to Hecke \(L\)-functions (the subject of his legendary PhD thesis), the Tate--Shafarevich group (the obstruction to the Hasse principle for elliptic curves), the Sato--Tate conjecture (the equidistribution of Frobenius traces), the Lubin--Tate theory of formal groups (which gives explicit local class field theory), the foundations of rigid analytic geometry (non-archimedean analytic geometry), the Tate curve (\(p\)-adic uniformization of elliptic curves), and the Tate conjecture (algebraic cycles and Galois representations). Each of these contributions is a chapter in the modern story of number theory.

\section{Biography}

John Torrence Tate Jr.\ was born on March 13, 1925, in Minneapolis, Minnesota, USA. His father, John Torrence Tate Sr., was a professor of physics at the University of Minnesota. The younger Tate developed an early interest in mathematics. He entered Harvard University in 1942, but his studies were interrupted by World War II, during which he served in the US Navy as a radio operator.

After the war, Tate completed his bachelor's degree at Harvard in 1946. He then moved to Princeton University for graduate study, where he wrote his legendary doctoral dissertation under the supervision of Emil Artin. Tate's 1950 PhD thesis, \textit{Fourier Analysis in Number Fields and Hecke's Zeta-Functions}, revolutionized the theory of \(L\)-functions by placing them in the framework of harmonic analysis on adeles. This work, though not published in full until years later, immediately established him as a major figure in number theory.

Tate spent the early part of his career at Princeton (1950--1953) and then moved to Harvard in 1953 as a Benjamin Peirce Instructor. He joined the faculty at Harvard and remained there until 1990, when he moved to the University of Texas at Austin. At Harvard, Tate supervised numerous PhD students, including Joe Buhler, Benedict Gross, Barry Mazur, and many others who would become leaders in arithmetic geometry.

Over the course of his career, Tate received almost every major honor in mathematics: the Wolf Prize in 2002/03 (shared with Mikio Sato), the Steele Prize for Lifetime Achievement in 1995, and the Abel Prize in 2010. He died on October 16, 2019, at the age of 94, leaving behind a body of work whose influence will be felt for generations.

\section{Tate Cohomology}

Group cohomology is a fundamental tool in number theory, providing the algebraic framework for class field theory, Galois representations, and the arithmetic of elliptic curves. For a finite group \(G\) and a \(G\)-module \(A\), the cohomology groups \(H^n(G,A)\) are defined for \(n \geq 0\) using the standard resolution. Tate discovered a remarkable extension of this theory to all integers \(n \in \Z\), unifying the homology and cohomology of finite groups into a single periodic, dimension-shifting theory.

\subsection{Definition via Complete Resolution}

Let \(G\) be a finite group and let \(A\) be a \(G\)-module. The group cohomology \(H^n(G,A)\) for \(n \geq 0\) is defined via the cochain complex of inhomogeneous cochains. The group homology \(H_n(G,A)\) for \(n \geq 0\) is defined via the bar resolution. Tate's key insight was to combine these into a \textit{complete resolution}: a \(\Z[G]\)-module resolution that is infinite in both directions.

\begin{definition}[Tate Cohomology Groups]
Let \(G\) be a finite group and \(A\) a \(G\)-module. The \textbf{Tate cohomology groups} \(\widehat{H}^n(G,A)\) are defined for all \(n \in \Z\) by
\[
\widehat{H}^n(G,A) =
\begin{cases}
H_n(G,A), & n \leq -2, \\[4pt]
\widehat{H}^{-1}(G,A) = \ker(N_G)/I_G A, & n = -1, \\[4pt]
\widehat{H}^{0}(G,A) = A^G / N_G(A), & n = 0, \\[4pt]
H^n(G,A), & n \geq 1,
\end{cases}
\]
where \(N_G: A \to A\) is the norm map \(N_G(a) = \sum_{g \in G} g \cdot a\), and \(I_G\) is the augmentation ideal of \(\Z[G]\).
\end{definition}

The norm map \(N_G\) and the augmentation ideal \(I_G\) sit at the heart of the theory. The group \(\widehat{H}^{-1}(G,A)\) measures the kernel of the norm modulo the augmentation ideal, while \(\widehat{H}^{0}(G,A)\) measures the \(G\)-invariants modulo the image of the norm. These two groups are connected by the \textbf{Tate--Nakayama theorem}, which is the cohomological engine of local class field theory.

\subsection{The Complete Resolution}

The complete resolution is constructed by splicing together the standard bar resolution (for negative degrees) and its dual (for positive degrees). Let \(P_\bullet \to \Z\) be a projective resolution of the trivial module \(\Z\) over \(\Z[G]\), and let \(\Z \to P^\bullet\) be an injective resolution. The complete resolution \(P_\bullet^{\text{Tate}}\) is a doubly-infinite complex
\[
\cdots \longrightarrow P_2 \longrightarrow P_1 \longrightarrow P_0 \longrightarrow P^0 \longrightarrow P^1 \longrightarrow P^2 \longrightarrow \cdots
\]
such that the cohomology of \(\Hom_{\Z[G]}(P_\bullet^{\text{Tate}}, A)\) computes \(\widehat{H}^n(G,A)\) for all \(n\).

\begin{theorem}[Properties of Tate Cohomology]
\begin{enumerate}
    \item \textbf{Long exact sequence:} For a short exact sequence \(0 \to A \to B \to C \to 0\) of \(G\)-modules, there is a long exact sequence
    \[
    \cdots \to \widehat{H}^{n}(G,A) \to \widehat{H}^{n}(G,B) \to \widehat{H}^{n}(G,C) \to \widehat{H}^{n+1}(G,A) \to \cdots
    \]
    infinite in both directions.
    \item \textbf{Periodicity:} If \(G\) is cyclic of order \(m\), then \(\widehat{H}^n(G,A) \cong \widehat{H}^{n+2}(G,A)\) for all \(n\).
    \item \textbf{Dimension shifting:} For any \(n\), \(\widehat{H}^{n+1}(G,A) \cong \widehat{H}^{n}(G, \Hom(\Z[G], A)/A)\).
    \item \textbf{Shapiro's lemma:} For a subgroup \(H \subseteq G\) and an \(H\)-module \(A\), \(\widehat{H}^n(H,A) \cong \widehat{H}^n(G, \Ind_H^G A)\).
\end{enumerate}
\end{theorem}

\subsection{The Tate--Nakayama Theorem and Class Field Theory}

The fundamental theorem of local class field theory can be expressed entirely in the language of Tate cohomology. For a finite Galois extension \(L/K\) of local fields with Galois group \(G = \Gal(L/K)\), the Brauer group \(\Br(L/K) = H^2(G, L^\times)\) is cyclic of order \([L:K]\), and the invariant map
\[
\inv_{L/K}: H^2(G, L^\times) \longrightarrow \frac{1}{[L:K]}\Z/\Z
\]
is an isomorphism. The Tate--Nakayama theorem provides the general cohomological framework:

\begin{theorem}[Tate--Nakayama]
Let \(G\) be a finite group acting on a module \(A\). Suppose that for every subgroup \(H \subseteq G\):
\begin{enumerate}
    \item \(H^1(H, A) = 0\), and
    \item \(H^2(H, A)\) is cyclic of order \(|H|\) with generator \(u_H\) whose restriction from \(G\) to \(H\) is compatible.
\end{enumerate}
Then for all \(n \in \Z\), cup product with the fundamental class \(u_G \in H^2(G, A)\) induces an isomorphism
\[
\cup u_G: \widehat{H}^n(G, \Z) \xrightarrow{\cong} \widehat{H}^{n+2}(G, A).
\]
\end{theorem}

This theorem is the cohomological heart of class field theory. When applied to \(A = L^\times\) (the multiplicative group of the completion of the maximal unramified extension), it yields the full structure of local class field theory: the existence theorem for norm subgroups, the reciprocity law, and the structure of the Weil group.

\subsection{The Herbrand Quotient}

For a finite group \(G\) and a \(G\)-module \(A\), the \textbf{Herbrand quotient} is defined when the cohomology groups are finite as
\[
h(G,A) = \frac{|\widehat{H}^0(G,A)|}{|\widehat{H}^{-1}(G,A)|} = \frac{|A^G/N_GA|}{|\ker N_G / I_G A|}.
\]
The Herbrand quotient is multiplicative in exact sequences and is invariant under induction. It plays a crucial role in the proofs of the Dirichlet unit theorem and the finiteness of the class group.

\section{Tate's Thesis}

Tate's 1950 Princeton PhD thesis, written under the supervision of Emil Artin, is one of the most celebrated dissertations in the history of mathematics. In it, Tate reinterpreted Hecke's theory of \(L\)-functions with Gr\"ossencharacters through the lens of harmonic analysis on the adeles, providing a unified and conceptual derivation of their analytic continuation and functional equation.

\subsection{Adeles and Ideles}

Let \(k\) be a number field with ring of integers \(\mathcal{O}_k\). For each place \(v\) of \(k\), let \(k_v\) be the completion. The \textbf{adele ring} of \(k\) is the restricted direct product
\[
\A_k = \prod_{v}' (k_v, \mathcal{O}_v),
\]
where \(\mathcal{O}_v\) is the ring of integers of \(k_v\) for non-archimedean \(v\). The restricted direct product means that for almost all \(v\), the component lies in \(\mathcal{O}_v\). The additive group \(\A_k\) is a locally compact topological ring. The \textbf{idele group} is the multiplicative analogue:
\[
\A_k^\times = \prod_{v}' (k_v^\times, \mathcal{O}_v^\times),
\]
which is a locally compact topological group. The multiplicative group \(k^\times\) embeds diagonally into \(\A_k^\times\) as a discrete cocompact subgroup (this is the content of the product formula and the Dirichlet unit theorem).

\subsection{Fourier Analysis on Adeles}

The key innovation of Tate's thesis was to apply Fourier analysis to the adeles. The additive group \(\A_k\) is self-dual: there exists a nontrivial additive character \(\psi: \A_k \to \C^\times\) such that the map
\[
x \mapsto \psi(x \cdot)
\]
identifies \(\A_k\) with its Pontryagin dual \(\widehat{\A}_k\). Tate constructed a Haar measure \(dx\) on \(\A_k\) normalized so that the Fourier inversion formula holds.

\begin{definition}[Tate's Zeta Integral]
Let \(\chi: \A_k^\times / k^\times \to \C^\times\) be a unitary Hecke character (a Gr\"ossencharacter of type \(A_0\)), and let \(f: \A_k \to \C\) be a Schwartz--Bruhat function (a function that is smooth and compactly supported at archimedean places, and locally constant with compact support at non-archimedean places). The \textbf{Tate zeta function} is
\[
\zeta(f, \chi, s) = \int_{\A_k^\times} f(a) \, \chi(a) \, |a|^s \, d^\times a,
\]
where \(d^\times a\) is the multiplicative Haar measure on \(\A_k^\times\) normalized so that \(\vol(\widehat{\mathcal{O}_k^\times}) = 1\) for the unit ideles.
\end{definition}

\subsection{Local Functional Equations}

At each place \(v\), Tate defined a local zeta integral
\[
\zeta_v(f_v, \chi_v, s) = \int_{k_v^\times} f_v(a_v) \, \chi_v(a_v) \, |a_v|_v^s \, d^\times a_v,
\]
where \(f_v\) is a Schwartz--Bruhat function on \(k_v\). The fundamental local result is:

\begin{theorem}[Local Functional Equation]
Let \(\psi_v\) be the local component of the global additive character \(\psi\). For each place \(v\), there exists a \textbf{local \(L\)-factor} \(L_v(\chi_v, s)\) and a \textbf{local epsilon factor} \(\varepsilon_v(\chi_v, s, \psi_v)\) such that for appropriate Schwartz--Bruhat functions \(f_v\),
\[
\frac{\zeta_v(\widehat{f}_v, \chi_v^{-1}, 1-s)}{L_v(\chi_v^{-1}, 1-s)}
= \varepsilon_v(\chi_v, s, \psi_v)\,
\frac{\zeta_v(f_v, \chi_v, s)}{L_v(\chi_v, s)},
\]
where \(\widehat{f}_v\) is the Fourier transform of \(f_v\) with respect to \(\psi_v\).
\end{theorem}

\subsection{Global Functional Equation}

Globally, Tate applied the Poisson summation formula on \(\A_k\):

\begin{theorem}[Tate's Thesis, Global Functional Equation]
Let \(k\) be a number field, \(\chi: \A_k^\times/k^\times \to \C^\times\) a Hecke character, and \(f: \A_k \to \C\) a Schwartz--Bruhat function. The completed \(L\)-function
\[
\Lambda(\chi, s) = \prod_v L_v(\chi_v, s) = \frac{\zeta(f, \chi, s)}{ \bigl( \int_{\A_k} f(x)\,dx \bigr)^{1/2 - s} }
\]
(up to a suitable normalization) satisfies the functional equation
\[
\Lambda(\chi, s) = \varepsilon(\chi, s) \, \Lambda(\chi^{-1}, 1-s),
\]
where \(\varepsilon(\chi, s) = \prod_v \varepsilon_v(\chi_v, s, \psi_v)\) is the global epsilon factor.
\end{theorem}

The global functional equation is derived by evaluating the integral
\[
\int_{\A_k^\times} f(a) \, \chi(a) \, |a|^s \, d^\times a
\]
in two ways: directly as the definition of \(\Lambda(\chi, s)\), and by applying Poisson summation on \(\A_k\) to the function \(f\). The key identity is
\[
\sum_{\xi \in k^\times} f(\xi a) = \frac{1}{|a|} \sum_{\eta \in k^\times} \widehat{f}(\eta / a),
\]
which is the adelic Poisson summation formula.

Tate's thesis unified the previously disparate theory of Hecke \(L\)-functions with Gr\"ossencharacters and Dirichlet \(L\)-functions, showing that all such \(L\)-functions arise from the same adelic construction. This framework later became the blueprint for the Langlands program, where automorphic representations are studied through their \(L\)-functions via harmonic analysis on adelic groups.

\section{The Tate--Shafarevich Group}

For elliptic curves over number fields, the Hasse principle --- the idea that a global solvability should follow from local solvability at all places --- can fail. The obstruction is measured by the Tate--Shafarevich group, denoted \(\Sha(E/k)\). This group is one of the central objects in arithmetic geometry and is intimately connected with the Birch and Swinnerton-Dyer conjecture.

\subsection{Definition}

Let \(E\) be an elliptic curve defined over a number field \(k\). For each place \(v\) of \(k\), let \(k_v\) be the completion and let \(G_v = \Gal(\overline{k}_v/k_v)\) be the local Galois group. Let \(G_k = \Gal(\overline{k}/k)\) be the absolute Galois group of \(k\).

\begin{definition}[Tate--Shafarevich Group]
The \textbf{Tate--Shafarevich group} of \(E\) over \(k\) is
\[
\Sha(E/k) = \ker\Bigl(\, H^1(G_k, E(\overline{k})) \longrightarrow \prod_v H^1(G_v, E(\overline{k}_v)) \,\Bigr).
\]
Equivalently, \(\Sha(E/k)\) consists of isomorphism classes of principal homogeneous spaces (torsors) for \(E\) over \(k\) that have \(k_v\)-rational points for every completion \(k_v\).
\end{definition}

An element of \(\Sha(E/k)\) is a curve \(C\) over \(k\) that is isomorphic to \(E\) over \(\overline{k}\) (so \(C\) is a twist of \(E\) by a cocycle) and that has points over every completion \(k_v\), but has no \(k\)-rational point. The group operation on \(\Sha(E/k)\) is given by the Baer sum of extensions, or equivalently by the group law on principal homogeneous spaces.

\subsection{Relation to the Selmer Group and the Mordell--Weil Group}

The Tate--Shafarevich group fits into the fundamental exact sequence relating the Mordell--Weil group \(E(k)\) to the Selmer groups. For any positive integer \(m\), the Kummer sequence gives
\[
0 \longrightarrow E(k)/mE(k) \longrightarrow H^1(G_k, E[m]) \longrightarrow H^1(G_k, E(\overline{k}))[m] \longrightarrow 0.
\]

The \(m\)-Selmer group is defined as the kernel
\[
\Sel^{(m)}(E/k) = \ker\Bigl(\, H^1(G_k, E[m]) \longrightarrow \prod_v H^1(G_v, E(\overline{k}_v)) \,\Bigr),
\]
and the Tate--Shafarevich group appears in the descent exact sequence
\[
0 \longrightarrow E(k)/mE(k) \longrightarrow \Sel^{(m)}(E/k) \longrightarrow \Sha(E/k)[m] \longrightarrow 0.
\]

This exact sequence is the primary tool for bounding the Mordell--Weil rank of an elliptic curve. The Selmer group is finite and effectively computable, while \(\Sha(E/k)\) is the principal obstruction to computing the Mordell--Weil group by descent.

\subsection{Conjectures}

The Tate--Shafarevich group is conjectured to be finite. This is one of the seven Clay Millennium Problems (the Birch and Swinnerton-Dyer conjecture) and one of the central open problems in arithmetic geometry.

\begin{conj}[Finiteness of \(\Sha\)]
For any elliptic curve \(E\) over a number field \(k\), the Tate--Shafarevich group \(\Sha(E/k)\) is finite.
\end{conj}

The finiteness of \(\Sha\) is known only in certain special cases:
\begin{itemize}
    \item For elliptic curves over \(\Q\) with analytic rank \(0\) or \(1\) (by the work of Kolyvagin and Rubin, building on the Gross--Zagier formula), the \(p\)-primary part of \(\Sha\) is finite.
    \item Over function fields, the finiteness of \(\Sha\) is known (by the work of Artin, Milne, and others).
    \item The full conjecture remains open for general elliptic curves over number fields.
\end{itemize}

\subsection{Relationship to the BSD Conjecture}

The Birch and Swinnerton-Dyer conjecture predicts that the rank of \(E(k)\) equals the order of vanishing of the Hasse--Weil \(L\)-function \(L(E/k, s)\) at \(s=1\), and that the leading coefficient of the Taylor expansion encodes deep arithmetic invariants including the order of \(\Sha(E/k)\). Specifically,
\[
\frac{L^{(r)}(E/k, 1)}{r!} = \frac{|\Sha(E/k)| \cdot \Omega_E \cdot R_E \cdot \prod_p c_p}{|E(k)_{\text{tors}}|^2},
\]
where \(\Omega_E\) is the real period, \(R_E\) is the regulator, \(c_p\) are the Tamagawa numbers, and \(E(k)_{\text{tors}}\) is the torsion subgroup.

\section{The Sato--Tate Conjecture}

The Sato--Tate conjecture, formulated jointly by Mikio Sato and John Tate in the 1960s, describes the distribution of the Frobenius traces of an elliptic curve. It is one of the most profound results in arithmetic geometry, finally proved by a team of mathematicians between 2006 and 2012.

\subsection{Elliptic Curves and Frobenius Traces}

Let \(E\) be an elliptic curve over \(\Q\). For a prime \(p\) of good reduction, let \(E_p\) denote the reduction of \(E\) modulo \(p\). The Hasse--Weil bound states that
\[
|a_p(E)| \leq 2\sqrt{p},
\]
where \(a_p(E) = p + 1 - |E_p(\F_p)|\) is the trace of Frobenius. The normalized trace is
\[
x_p = \frac{a_p(E)}{2\sqrt{p}} \in [-1, 1].
\]

Sato, through extensive numerical computation on the IBM 7090 at the University of Tokyo, observed that the distribution of the \(x_p\) for curves without complex multiplication seemed to follow a specific semicircular law. Tate independently arrived at the same conjecture from a theoretical perspective.

\begin{conj}[Sato--Tate Conjecture, 1960s]
Let \(E\) be an elliptic curve over \(\Q\) without complex multiplication. Then as \(p\) varies over primes of good reduction, the normalized Frobenius traces \(x_p = a_p(E)/2\sqrt{p}\) are equidistributed in \([-1, 1]\) with respect to the measure
\[
\mu_{\ST}(x) = \frac{2}{\pi} \, \sqrt{1 - x^2} \; dx.
\]
\end{conj}

Equivalently, for any continuous function \(f: [-1, 1] \to \C\),
\[
\lim_{X \to \infty} \frac{1}{\pi(X)} \sum_{p \leq X} f(x_p) = \int_{-1}^{1} f(x) \, \frac{2}{\pi} \sqrt{1-x^2} \, dx,
\]
where \(\pi(X)\) is the prime counting function.

The measure \(\mu_{\ST}\) is the Sato--Tate measure, which is the pushforward of the Haar measure on \(\SU(2)\) under the trace map. The Sato--Tate group in this case is \(\SU(2)\).

\subsection{The Sato--Tate Theorem}

The conjecture was proved through the collective effort of Clozel, Harris, Shepherd-Barron, and Taylor (2006--2012), building on earlier work of Wiles and others on modularity lifting theorems.

\begin{theorem}[Sato--Tate Theorem, 2006--2012]
Let \(E\) be an elliptic curve over a totally real number field \(F\) without complex multiplication. Then the normalized Frobenius traces \(a_{\mathfrak{p}}(E)/2\sqrt{N(\mathfrak{p})}\) are equidistributed with respect to the Sato--Tate measure. For curves with complex multiplication, the distribution is different, corresponding to the Haar measure on the maximal torus of \(\SU(2)\).
\end{theorem}

The proof proceeds by establishing the potential automorphy of the symmetric power \(L\)-functions of \(E\). For each \(n \geq 1\), the symmetric power \(L\)-function \(L(\Sym^n E, s)\) is shown to be meromorphic with a functional equation and to be analytic in a region containing the critical strip, and the absence of poles or zeros on the edge of the critical strip implies the equidistribution theorem via the Weyl criterion.

\subsection{Generalizations}

The Sato--Tate conjecture has been generalized far beyond elliptic curves. For any smooth projective variety \(X\) over a number field, the Frobenius eigenvalues at each prime produce a statistical distribution governed by a compact Lie group --- the Sato--Tate group of \(X\). The generalized Sato--Tate conjecture (formulated by Serre, Katz, and others) predicts:

\begin{conj}[Generalized Sato--Tate Conjecture]
Let \(X\) be a smooth projective variety over a number field \(k\). For each prime \(v\) of good reduction, consider the normalized Frobenius eigenvalues \(\{\alpha_{i,v}\}\) acting on the \(l\)-adic cohomology \(H^i_{\text{\'et}}(X_{\overline{k}}, \Q_l)\). These eigenvalues are equidistributed in the unitary group with respect to the Haar measure of the Sato--Tate group \(\ST(X)\), which is the Zariski closure of the image of the Galois group in the automorphism group of the cohomology.
\end{conj}

This conjecture has been proved for abelian varieties of GL(2)-type by the work of Fite, Kedlaya, Rotger, and Sutherland, and for many specific families of varieties through the ongoing development of potential automorphy theorems.

\section{Formal Groups and Lubin--Tate Theory}

The theory of formal groups, developed by Tate in collaboration with Jonathan Lubin, provides explicit constructions for local class field theory. The Lubin--Tate theory gives a remarkably concrete description of the maximal abelian extension of a local field, using the arithmetic of formal group laws.

\subsection{Formal Group Laws}

A (commutative) one-dimensional formal group law over a ring \(R\) is a power series \(F(X,Y) \in R[[X,Y]]\) satisfying
\begin{align*}
F(X,0) &= X, \quad F(0,Y) = Y,\\
F(F(X,Y),Z) &= F(X,F(Y,Z)),\\
F(X,Y) &= F(Y,X).
\end{align*}
The formal group law endows the set of power series with zero constant term with a group structure via \(X +_F Y = F(X,Y)\).

\begin{example}[The Additive and Multiplicative Groups]
The additive formal group is \(\widehat{\G}_a(X,Y) = X + Y\). The multiplicative formal group is \(\widehat{\G}_m(X,Y) = X + Y + XY = (1+X)(1+Y) - 1\).
\end{example}

\subsection{The Lubin--Tate Construction}

Let \(k\) be a finite extension of \(\Q_p\) with residue field \(\F_q\). Let \(\pi\) be a uniformizer of \(k\). The Lubin--Tate formal group is associated to the Frobenius endomorphism.

\begin{definition}[Lubin--Tate Formal Group]
Let \(f(X) = \pi X + X^q\). There exists a unique formal group law \(F_f(X,Y)\) over \(\mathcal{O}_k\) such that
\[
f(F_f(X,Y)) = F_f(f(X), f(Y)),
\]
and such that \(f\) is an endomorphism of \(F_f\). The formal group \(F_f\) is called the \textbf{Lubin--Tate formal group} associated to \((k,\pi)\).
\end{definition}

\begin{theorem}[Lubin--Tate, 1965]
Let \(k_\pi\) be the field obtained by adjoining to \(k\) the torsion points of the Lubin--Tate formal group \(F_f\). Then:
\begin{enumerate}
    \item \(k_\pi\) is the maximal abelian extension of \(k\).
    \item The Galois group \(\Gal(k_\pi/k)\) is canonically isomorphic to \(\widehat{\mathcal{O}_k^\times}\), the profinite completion of \(\mathcal{O}_k^\times\).
    \item The extension \(k_\pi\) is totally ramified, and its residue field is the algebraic closure of \(\F_q\).
\end{enumerate}
\end{theorem}

The key object is the \textbf{formal module} \(\mathfrak{m} = (\mathfrak{m}_{\overline{k}}, +_{F_f})\), the set of maximal ideal elements equipped with the group law from \(F_f\). The torsion points \(F_f[\pi^n] = \{x \in \mathfrak{m} : f^n(x) = 0\}\) generate the Lubin--Tate tower of extensions.

\subsection{The Local Reciprocity Map}

The Lubin--Tate theory provides an explicit construction of the local reciprocity map (the Artin map). For a finite extension \(L/k\), the norm residue symbol is realized through the action of the Galois group on the torsion points of the formal group:

\[
(\cdot, L/k): k^\times \longrightarrow \Gal(L/k)^{\ab}.
\]

Via the Lubin--Tate construction, this map is given by the formula: for \(a \in \mathcal{O}_k^\times\), the corresponding element of \(\Gal(k_\pi/k)\) acts on the torsion points \(F_f[\pi^n]\) by multiplication by the image of \(a\) in \(\mathcal{O}_k^\times/(1 + \pi^n \mathcal{O}_k)\), made explicit through the formal group exponential. This gives the cleanest and most concrete construction of local class field theory.

\subsection{Explicit Formulas}

The Lubin--Tate theory also yields explicit formulas for the Hilbert symbol (the reciprocity law for the \(p\)-power cyclotomic extensions). For a local field \(k\) containing the \(p^m\)-th roots of unity, the \(p^m\)-th Hilbert symbol
\[
(\cdot, \cdot)_{p^m}: k^\times \times k^\times \longrightarrow \mu_{p^m}
\]
can be expressed in terms of the formal group logarithm and the Artin--Hasse exponential series. These formulas have profound applications to the explicit reciprocity laws in Iwasawa theory.

\section{Rigid Analytic Geometry}

The foundations of non-archimedean analytic geometry were laid by Tate in a series of seminars and his 1971 paper "Rigid Analytic Spaces." Prior to Tate's work, the theory of analytic functions over non-archimedean fields was severely limited by the total disconnectedness of the topology: because the absolute value is ultrametric, any two disjoint open sets are separated, making the usual complex-analytic approach (based on the identity theorem and analytic continuation) fail. Tate's revolutionary idea was to replace the naive topology with a Grothendieck topology (now called the rigid topology), allowing a full-fledged theory of analytic spaces.

\subsection{Tate Algebras}

Let \(k\) be a complete non-archimedean field with valuation \(|\cdot|\) and residue field \(\tilde{k}\).

\begin{definition}[Tate Algebra]
The \textbf{Tate algebra} in \(n\) variables over \(k\) is
\[
T_n = k\langle \xi_1, \dots, \xi_n \rangle = \Bigl\{ \sum_{|\alpha| \geq 0} a_\alpha \xi^\alpha \in k[[\xi_1,\dots,\xi_n]] : |a_\alpha| \to 0 \text{ as } |\alpha| \to \infty \Bigr\}.
\]
These are the power series that converge on the unit polydisk \(\{(x_1,\dots,x_n) \in k^n : |x_i| \leq 1\}\).
\end{definition}

Tate algebras are noetherian and satisfy the Weierstrass preparation and division theorems, making them the non-archimedean analogues of the ring of convergent power series over \(\C\).

\subsection{Affinoid Spaces}

Let \(A = T_n / I\) be a quotient of a Tate algebra by an ideal \(I\). The \textbf{affinoid space} associated to \(A\) is the set \(\Sp(A)\) of maximal ideals of \(A\) (equivalently, bounded multiplicative seminorms on \(A\)). The topology on \(\Sp(A)\) is the \textbf{rigid analytic topology}, a Grothendieck topology generated by rational subdomains.

\begin{definition}[Rational Subdomain]
A subset \(U \subseteq \Sp(A)\) is a \textbf{rational subdomain} if there exist elements \(f_0, \dots, f_m \in A\) generating the unit ideal such that
\[
U = \{x \in \Sp(A) : |f_i(x)| \leq |f_0(x)| \text{ for all } i = 1, \dots, m\}.
\]
\end{definition}

The rational subdomains form a basis for the \textbf{weak topology} on \(\Sp(A)\), and the associated Grothendieck topology (with admissible coverings given by finite families of rational subdomains whose images in the valuation spectrum cover the space) is called the rigid topology.

\begin{theorem}[Tate's Acyclicity Theorem, 1971]
For any affinoid space \(\Sp(A)\) and any admissible covering \(\{U_i\}\) by rational subdomains, the \v{C}ech cohomology of the structure sheaf \(\mathcal{O}\) vanishes:
\[
\check{H}^q(\{U_i\}, \mathcal{O}) = 0 \quad \text{for all } q \geq 1.
\]
Consequently, the structure sheaf is a sheaf on the rigid topology.
\end{theorem}

\subsection{Rigid Analytic Spaces}

Rigid analytic spaces are obtained by gluing affinoid spaces along admissible open subspaces, in perfect analogy with the glued construction of schemes in algebraic geometry. The resulting category of rigid analytic spaces admits:
\begin{itemize}
    \item A theory of coherent sheaves and their cohomology (Kiehl's theorems).
    \item The GAGA theorem: for a proper scheme \(X\) over \(k\), the analytification functor \(X \mapsto X^{\an}\) induces an equivalence of categories between coherent algebraic sheaves on \(X\) and coherent analytic sheaves on \(X^{\an}\).
    \item The theory of etale cohomology for rigid spaces, with comparison theorems to algebraic etale cohomology.
\end{itemize}

The foundations laid by Tate were later greatly expanded by Raynaud, Berthelot, Berkovich, Huber, and others, leading to the modern theories of Berkovich spaces, Huber's adic spaces, and the theory of perfectoid spaces developed by Scholze.

\section{The Tate Curve}

The Tate curve is the \(p\)-adic uniformization of elliptic curves, giving a construction of elliptic curves over \(p\)-adic fields as quotients of the multiplicative group by a lattice. It is the non-archimedean analogue of the classical uniformization of elliptic curves over \(\C\) as \(\C/\Lambda\).

\subsection{Classical Uniformization}

Over the complex numbers, every elliptic curve \(E\) is isomorphic to a torus \(\C/(\Z + \tau\Z)\) for some \(\tau \in \mathfrak{h}\) (the upper half-plane). The corresponding Weierstrass parametrization gives
\[
\wp(z; \Lambda) = \frac{1}{z^2} + \sum_{\omega \in \Lambda \setminus \{0\}} \left( \frac{1}{(z-\omega)^2} - \frac{1}{\omega^2} \right),
\]
and the map \(z \mapsto (\wp(z), \wp'(z))\) identifies \(\C/\Lambda\) with the elliptic curve \(y^2 = 4x^3 - g_2 x - g_3\).

\subsection{The Tate Curve Over \(p\)-adic Fields}

Over a \(p\)-adic field \(k\), the situation is radically different. An elliptic curve with bad reduction (specifically, multiplicative reduction) admits a uniformization by the multiplicative group \(\G_m\).

\begin{definition}[Tate Curve]
For a parameter \(q \in k^\times\) with \(|q| < 1\), the \textbf{Tate curve} \(E_q\) is the elliptic curve over \(k\) defined by the Weierstrass equation
\[
y^2 + xy = x^3 + a_4(q) x + a_6(q),
\]
where
\[
a_4(q) = -5 \sum_{n \geq 1} \frac{n^3 q^n}{1-q^n}, \qquad
a_6(q) = -\frac{1}{12} \sum_{n \geq 1} \frac{(7n^5 + 5n^3)q^n}{1-q^n}.
\]
The \(j\)-invariant of \(E_q\) is given by the power series
\[
j(q) = \frac{1}{q} + 744 + 196884 q + 21493760 q^2 + \cdots,
\]
which recovers the classical \(j\)-function.
\end{definition}

\begin{theorem}[Tate, 1960s]
For any \(q \in k^\times\) with \(|q| < 1\), the Tate curve \(E_q\) satisfies:
\begin{enumerate}
    \item The map \(\phi: \overline{k}^\times / q^\Z \to E_q(\overline{k})\) given by
    \[
    \phi(u) = \left( \sum_{n \in \Z} \frac{q^n u}{(1-q^n u)^2} - 2 \sum_{n \geq 1} \frac{n q^n}{1-q^n}, \;
    \sum_{n \in \Z} \frac{q^{2n} u^2}{(1 - q^n u)^3} + \sum_{n \geq 1} \frac{n q^n}{1-q^n} \right)
    \]
    is an analytic isomorphism of rigid analytic groups.
    \item The curve \(E_q\) has split multiplicative reduction: the special fiber of the N\'eron model is a nodal cubic.
    \item The \(q\)-invariant \(j(q)\) is the \(j\)-invariant of \(E_q\), and the assignment \(q \mapsto j(q)\) gives a bijection between isomorphism classes of Tate curves and elements of \(k^\times\) with \(|j| > 1\).
\end{enumerate}
\end{theorem}

The Tate curve provides the \(p\)-adic analogue of the exponential map \(z \mapsto e^{2\pi i z}\) that uniformizes complex elliptic curves. The uniformization \(\overline{k}^\times / q^\Z \cong E_q(\overline{k})\) is given by the \(q\)-expansion of the Weierstrass \(\wp\)-function, with the theta functions taking the place of trigonometric functions.

\subsection{Applications}

The Tate curve has numerous applications in arithmetic geometry:
\begin{itemize}
    \item \textbf{Tamagawa numbers:} The Tate curve provides an explicit formula for the Tamagawa number \(c_v\) of an elliptic curve with multiplicative reduction at a place \(v\): \(c_v = \ord_v(q)\).
    \item \textbf{Maps to the Tate curve:} For an elliptic curve \(E\) with multiplicative reduction, the Tate uniformization gives an explicit description of the connected component of the N\'eron model and the group of components.
    \item \textbf{Hecke theory:} The \(q\)-expansion of modular forms at cusps is the limit of the Tate uniformization as \(q \to 0\), connecting the arithmetic theory of modular forms with the geometry of elliptic curves.
    \item \textbf{Local heights:} The Tate curve yields explicit formulas for N\'eron's local height pairing, essential for the Gross--Zagier formula.
\end{itemize}

\section{The Tate Conjecture}

The Tate conjecture, formulated in the 1960s, is a sweeping statement linking algebraic cycles on algebraic varieties to Galois representations. It is one of the central problems in arithmetic geometry and is closely related to the Hodge conjecture.

\subsection{Statement}

Let \(X\) be a smooth projective variety of dimension \(d\) over a finitely generated field \(k\). Choose a prime \(l \neq \operatorname{char}(k)\) and consider the etale cohomology groups \(H^r_{\text{\'et}}(X_{\overline{k}}, \Q_l(r))\). These are representations of the absolute Galois group \(G_k = \Gal(\overline{k}/k)\).

For an algebraic cycle \(Z \subset X\) of codimension \(r\) defined over \(k\), the cycle class map associates a cohomology class
\[
cl(Z) \in H^{2r}_{\text{\'et}}(X_{\overline{k}}, \Q_l(r)),
\]
which is fixed by the action of \(G_k\) (because \(Z\) is defined over \(k\)).

\begin{conj}[Tate Conjecture, 1963]
Let \(X\) be a smooth projective variety over a finitely generated field \(k\). For any prime \(l \neq \operatorname{char}(k)\), the cycle class map
\[
\Pic^r(X) \longrightarrow H^{2r}_{\text{\'et}}(X_{\overline{k}}, \Q_l(r))^{G_k}
\]
is surjective when tensored with \(\Q_l\). More precisely, the \(\Q_l\)-vector space spanned by classes of algebraic cycles of codimension \(r\) (defined over \(k\)) coincides with the subspace of \(H^{2r}_{\text{\'et}}(X_{\overline{k}}, \Q_l(r))\) that is fixed by the Galois group \(G_k\).
\end{conj}

In the simplest case \(r = 1\), the Tate conjecture asserts that the Galois-invariant part of \(H^2_{\text{\'et}}(X_{\overline{k}}, \Q_l(1))\) is spanned by the classes of divisors. This case is equivalent to the finiteness of the Brauer group of \(X\) (the Tate conjecture for divisors), and it follows from the Kummer sequence and the theory of N\'eron--Severi groups once the algebraic closure of the base field is finite.

\begin{definition}[Tate Conjecture, Divisor Case]
For a smooth projective variety \(X\) over a finitely generated field \(k\), the Tate conjecture for divisors (\(r=1\)) asserts:
\[
\NS(X_{\overline{k}}) \otimes \Q_l \cong H^2_{\text{\'et}}(X_{\overline{k}}, \Q_l(1))^{G_k},
\]
where \(\NS\) is the N\'eron--Severi group.
\end{definition}

\subsection{Relation to the Hodge Conjecture}

The Tate conjecture is the Galois-theoretic analogue of the Hodge conjecture. Over \(\C\), the comparison isomorphism between etale cohomology and singular cohomology identifies the Galois-invariant subspace with the subspace of Hodge classes. The Hodge conjecture predicts that every Hodge class on a smooth projective complex variety comes from a rational linear combination of algebraic cycles.

The analogy is:
\[
\begin{array}{ccc}
\text{Tate conjecture (Galois invariance)} & \longleftrightarrow & \text{Hodge conjecture (Hodge weights)} \\
\text{Gal}_{\overline{k}/k}\text{-invariant classes} & & (p,p)\text{-classes under Hodge decomposition}
\end{array}
\]

\subsection{Cases Proved}

The Tate conjecture has been proved in the following cases:
\begin{itemize}
    \item \textbf{Divisors on abelian varieties} (Tate, 1966): For abelian varieties over finite fields, the Tate conjecture for divisors holds. This was generalized by Zarhin and Faltings to abelian varieties over number fields.
    \item \textbf{Abelian varieties over finite fields} (Tate): The full Tate conjecture (all codimensions) for abelian varieties over finite fields was proved by Tate for the product of two curves and by Zarhin, Mori, and others in full generality.
    \item \textbf{K3 surfaces over finite fields} (Nygaard, Ogus, Charles, Madapusi Pera): The Tate conjecture for K3 surfaces over finite fields was proved in the 2010s, building on the work of Nygaard and Ogus for ordinary K3 surfaces.
    \item \textbf{Products of curves} (Tate, 1994): For products \(C_1 \times C_2\) of curves over a finite field, the Tate conjecture for codimension 2 was proved using the theory of correspondences and the Riemann hypothesis for curves.
\end{itemize}

\subsection{Consequences and Relationships}

The Tate conjecture is deeply connected to many other fundamental conjectures in arithmetic geometry:
\begin{itemize}
    \item \textbf{Finiteness of Brauer groups:} The Tate conjecture for divisors implies the finiteness of the Brauer group \(\Br(X)\) (and the Tate conjecture in higher codimensions implies the finiteness of higher Brauer groups).
    \item \textbf{The Birch and Swinnerton-Dyer conjecture:} For an elliptic curve \(E\), the Tate conjecture for the product \(E \times E\) is equivalent to the BSD conjecture for \(E\) (by the work of Tate and Beilinson).
    \item \textbf{Generalized Hodge conjecture:} Over number fields, the Tate conjecture for all varieties implies the Hodge conjecture for all varieties defined over number fields (by the comparison isomorphism and the standard conjectures).
\end{itemize}

\section{Impact and Legacy}

\subsection{Number Theory and Arithmetic Geometry}

John Tate's work transformed number theory from a subject rooted in classical analytic methods to a deep and interconnected edifice built on algebraic geometry, representation theory, and homological algebra. His creation of Tate cohomology provided the language for class field theory, and his thesis recast the theory of \(L\)-functions on a firm adelic foundation that directly anticipates the Langlands program.

\subsection{The Tate--Shafarevich Group and the BSD Conjecture}

The Tate--Shafarevich group \(\Sha(E/k)\) is one of the central objects of arithmetic geometry. The conjecture that it is finite is equivalent to the second part of the Birch and Swinnerton-Dyer conjecture, which is one of the Clay Millennium Problems. The study of \(\Sha\) and its relationship to Selmer groups, descent theory, and \(p\)-adic methods remains one of the most active areas of research.

\subsection{The Sato--Tate Conjecture}

The proof of the Sato--Tate conjecture by Clozel, Harris, Shepherd-Barron, and Taylor represents one of the most spectacular achievements of 21st-century number theory. It confirmed a deep intuition of Tate and Sato about the random behavior of Frobenius traces and opened the door to a vast generalization: the study of Sato--Tate groups for arbitrary motives, now a thriving field of research.

\subsection{Rigid Analytic Geometry and the Tate Curve}

Tate's rigid analytic geometry laid the foundations for modern non-archimedean geometry. The theory of the Tate curve is essential for understanding elliptic curves with multiplicative reduction. These ideas were later developed by Raynaud (formal geometry), Berkovich (Berkovich spaces), Huber (adic spaces), and most recently Scholze (perfectoid spaces and the theory of diamonds), which underpin the revolutionary progress in \(p\)-adic Hodge theory and the Langlands program.

\subsection{The Tate Conjecture}

The Tate conjecture remains one of the grand open problems of arithmetic geometry. Its proof for divisors on abelian varieties and K3 surfaces over finite fields represents significant progress, but the full conjecture is far from resolved. The conjecture is a central pillar of the web of conjectures (together with the Hodge conjecture, Beilinson's conjectures, and the standard conjectures) that define the landscape of modern arithmetic geometry.

\section{Timeline}

\begin{tabular}{|l|l|}
\hline
\textbf{Year} & \textbf{Event} \\ \hline
1925 & Born on March 13 in Minneapolis, Minnesota, USA \\ \hline
1942 & Entered Harvard University; studies interrupted by WWII \\ \hline
1946 & A.B.\ in Mathematics, Harvard University \\ \hline
1950 & Ph.D.\ in Mathematics, Princeton University (advisor: Emil Artin); Tate's thesis \\ \hline
1950--1953 & Member, Institute for Advanced Study, Princeton \\ \hline
1953 & Benjamin Peirce Instructor, Harvard University \\ \hline
1954 & Joined the faculty at Harvard University \\ \hline
1950s & Development of Tate cohomology and modification of group cohomology \\ \hline
1950s & Tate's thesis on Fourier analysis in number fields \\ \hline
1960s & Discovery of the Tate--Shafarevich group \\ \hline
1960s & Formulation of the Tate conjecture on algebraic cycles \\ \hline
1960s & Tate curve: \(p\)-adic uniformization of elliptic curves \\ \hline
1960s & Sato--Tate conjecture formulated jointly with Mikio Sato \\ \hline
1960s & Lubin--Tate theory of formal groups (with Jonathan Lubin) \\ \hline
1971 & Rigid analytic geometry foundations published \\ \hline
1990 & Moved to the University of Texas at Austin \\ \hline
1995 & Awarded the Steele Prize for Lifetime Achievement \\ \hline
2002/03 & Awarded the Wolf Prize in Mathematics (shared with Mikio Sato) \\ \hline
2006--12 & Sato--Tate conjecture proved (Clozel--Harris--Shepherd-Barron--Taylor) \\ \hline
2010 & Awarded the Abel Prize from the Norwegian Academy of Science and Letters \\ \hline
2019 & Died on October 16 in Lexington, Massachusetts, USA \\ \hline
\end{tabular}

\section*{Frequently Asked Questions}
\addcontentsline{toc}{section}{Frequently Asked Questions}

\subsection*{Q1: What is Tate cohomology and why is it important?}

Tate cohomology extends ordinary group cohomology \(H^n(G,A)\) (defined for \(n \geq 0\)) to all integers \(n \in \Z\), unifying group homology and cohomology. It provides the cohomological framework for class field theory: the Tate--Nakayama theorem, which relates the Brauer group to the formation of class formations, is the engine behind the reciprocity laws of local and global class field theory.

\subsection*{Q2: What was the main contribution of Tate's thesis?}

Tate's thesis reinterpreted Hecke \(L\)-functions with Gr\"ossencharacters through harmonic analysis on the adeles. By constructing a local zeta integral and applying Poisson summation on the adele ring, Tate gave a clean, unified proof of the analytic continuation and functional equation of all such \(L\)-functions. This adelic framework directly anticipates the Langlands program, where automorphic representations replace Gr\"ossencharacters.

\subsection*{Q3: What does the Tate--Shafarevich group measure?}

The Tate--Shafarevich group \(\Sha(E/k)\) measures the failure of the Hasse principle for an elliptic curve \(E\) over a number field \(k\). Its elements are curves (principal homogeneous spaces) that have rational points at every completion \(k_v\) but no global rational points. The finiteness of \(\Sha\) is equivalent to the second part of the Birch and Swinnerton-Dyer conjecture.

\subsection*{Q4: What does the Sato--Tate conjecture say and why is it important?}

The Sato--Tate conjecture states that the normalized Frobenius traces \(a_p(E)/2\sqrt{p}\) of an elliptic curve without complex multiplication are distributed according to a semicircular law (the Sato--Tate measure). It was proved between 2006 and 2012. The conjecture is important because it reveals the deep statistical behavior of arithmetic data and has been generalized to arbitrary motives.

\subsection*{Q5: What is the Lubin--Tate theory of formal groups?}

Lubin--Tate theory provides an explicit construction of the maximal abelian extension of a local field using torsion points of a carefully chosen formal group law. This gives the most concrete construction of local class field theory and yields explicit formulas for the reciprocity map and the Hilbert symbol.

\subsection*{Q6: What is rigid analytic geometry?}

Rigid analytic geometry is Tate's theory of non-archimedean analytic geometry. Because non-archimedean fields are totally disconnected, the naive analytic topology is useless. Tate replaced it with a Grothendieck topology (the rigid topology) whose admissible coverings are required to be finite, leading to a rich theory of analytic spaces over non-archimedean fields.

\subsection*{Q7: What is the Tate curve?}

The Tate curve is the \(p\)-adic uniformization of an elliptic curve with multiplicative reduction: \(E_q \cong \overline{k}^\times / q^\Z\). It is the non-archimedean analogue of the classical uniformization \(\C/\Lambda\). The parameter \(q\) gives the \(j\)-invariant via \(j(q) = 1/q + 744 + \cdots\).

\subsection*{Q8: What is the current status of the Tate conjecture?}

The Tate conjecture is proved for divisors on abelian varieties (Tate, Zarhin, Faltings) and for K3 surfaces over finite fields. It is known for several other special classes of varieties but is wide open in general. It is one of the central open problems in arithmetic geometry, closely related to the Hodge conjecture and the BSD conjecture.

\subsection*{Q9: How did Tate's work influence the Langlands program?}

Tate's adelic framework in his thesis was a direct precursor to Langlands' re-interpretation of automorphic forms as functions on \(\GL_n(\A_k)\) and the Langlands philosophy of associating Galois representations to automorphic representations. The symmetric power \(L\)-functions used in the proof of the Sato--Tate conjecture are a key instance of the Langlands functoriality principle.

\section*{Exercises}
\addcontentsline{toc}{section}{Exercises}

\begin{enumerate}
    \item Let \(G\) be a finite cyclic group of order \(m\) with generator \(g\). Show that \(\widehat{H}^0(G, \Z) \cong \Z/m\Z\) and \(\widehat{H}^{-1}(G, \Z) = 0\). Compute \(\widehat{H}^n(G, \Z)\) for all \(n\).
    \item Let \(A\) be a \(G\)-module. Prove that \(\widehat{H}^0(G, A) = A^G / N_G A\) and that \(\widehat{H}^{-1}(G, A) = \ker N_G / I_G A\).
    \item Show that the Herbrand quotient is multiplicative: for a short exact sequence \(0 \to A \to B \to C \to 0\) of \(G\)-modules, \(h(G,B) = h(G,A) \cdot h(G,C)\) whenever all three are defined.
    \item For the multiplicative formal group \(\widehat{\G}_m(X,Y) = X + Y + XY\), find an explicit isomorphism to the additive formal group \(\widehat{\G}_a\) over a ring containing \(\Q\).
    \item Using the Lubin--Tate construction, prove that for a uniformizer \(\pi\) of \(\Q_p\), the abelian extension \(\Q_p(\mu_{p^\infty})\) is generated by the torsion points of the multiplicative formal group.
    \item Prove that for the Tate curve \(E_q\), the map \(\phi: \overline{k}^\times / q^\Z \to E_q(\overline{k})\) is a bijection by showing that the Weierstrass functions defined by the series converge for all \(u \notin q^\Z\).
    \item Show that \(j(q) = q^{-1} + 744 + 196884 q + \cdots\) and verify that the coefficients are integers. (The coefficient 196884 has a deep relationship with the Monster group and monstrous moonshine.)
    \item Let \(E\) be an elliptic curve over \(\Q\) with CM by an imaginary quadratic order. Describe the Sato--Tate distribution of \(E\) and compare it with the Sato--Tate measure for non-CM curves.
    \item Prove that for an elliptic curve \(E\) without CM, the Sato--Tate group is \(\SU(2)\). Show that the trace of the standard representation of \(\SU(2)\) gives the measure \(\frac{2}{\pi}\sqrt{1-x^2}\,dx\).
    \item Let \(X\) be a smooth projective variety over a finite field \(k\). Show that the Tate conjecture for divisors on \(X\) is equivalent to the finiteness of the Brauer group \(\Br(X)\).
    \item For an elliptic curve \(E\) over a number field \(k\), show that \(\Sha(E/k)[m]\) fits into the exact sequence
    \[
    0 \to E(k)/mE(k) \to \Sel^{(m)}(E/k) \to \Sha(E/k)[m] \to 0.
    \]
    \item Let \(f\) be a Schwartz--Bruhat function on the adeles of \(\Q\). Show that the Poisson summation formula on \(\A_\Q\) gives
    \[
    \sum_{\xi \in \Q} f(\xi) = \sum_{\eta \in \Q} \widehat{f}(\eta).
    \]
    \item Verify that the local epsilon factors \(\varepsilon_v(\chi_v, s, \psi_v)\) in Tate's thesis satisfy the product formula \(\prod_v \varepsilon_v(\chi_v, s, \psi_v) = 1\), giving consistency with the global functional equation.
    \item Let \(T_n = k\langle \xi_1, \dots, \xi_n \rangle\) be a Tate algebra. Prove that \(T_n\) is a Noetherian ring and that its Jacobson radical is the set of power series with constant term in the maximal ideal of \(k\).
    \item Show that for \(X = \Sp(T_n)\), the rational subdomains form a basis for a Grothendieck topology and that the structure sheaf \(\mathcal{O}\) is a sheaf with respect to admissible coverings.
    \item Let \(E\) be an elliptic curve with split multiplicative reduction at a prime \(p\). Use the Tate uniformization to compute the Tamagawa number \(c_p\) in terms of the valuation of the Tate parameter \(q\).
    \item Describe the relationship between the Tate conjecture for \(X \times X\) and the BSD conjecture for an elliptic curve \(X\).
\end{enumerate}

\section*{Glossary}
\addcontentsline{toc}{section}{Glossary}

\begin{description}
    \item[Adele ring] The restricted direct product of all completions of a number field, a locally compact topological ring.
    \item[Affinoid space] The space of maximal ideals of a Tate algebra, the basic building block of rigid analytic geometry.
    \item[Class field theory] The theory that describes the abelian extensions of a local or global field in terms of the multiplicative group or the idele class group.
    \item[Complete resolution] A doubly-infinite projective resolution used to define Tate cohomology for all integer degrees.
    \item[Formal group law] A power series \(F(X,Y)\) defining a group structure on the set of elements with zero constant term.
    \item[Gr\"ossencharacter] A continuous homomorphism \(\A_k^\times/k^\times \to \C^\times\), a Hecke character in the sense of Tate's thesis.
    \item[Hasse principle] The principle that a global Diophantine problem has a solution if and only if it has solutions at every local completion.
    \item[Ideles] The multiplicative group of the adele ring, a locally compact topological group.
    \item[Lubin--Tate formal group] A formal group over the ring of integers of a local field whose torsion points generate the maximal abelian extension.
    \item[Sato--Tate conjecture] The conjecture, now theorem, that normalized Frobenius traces are equidistributed with respect to a semicircular law.
    \item[Sato--Tate group] A compact Lie group whose Haar measure governs the distribution of Frobenius eigenvalues.
    \item[Tate algebra] A convergent power series ring over a non-archimedean field, the basic object of rigid analytic geometry.
    \item[Tate cohomology] A unified cohomology theory for finite groups defined for all integer degrees \(n \in \Z\).
    \item[Tate conjecture] The conjecture that Galois-invariant etale cohomology classes are spanned by classes of algebraic cycles.
    \item[Tate curve] The \(p\)-adic uniformization of an elliptic curve with multiplicative reduction as \(\overline{k}^\times / q^\Z\).
    \item[Tate--Nakayama theorem] A cohomological theorem that provides the foundation for class formations and the reciprocity isomorphism.
    \item[Tate--Shafarevich group] The group \(\Sha(E/k)\) classifying principal homogeneous spaces with local but not global rational points.
    \item[Tamagawa number] The volume of the adelic points of an algebraic group with respect to the Tamagawa measure, given by an Euler product of local factors.
    \item[Theta function] A power series \(\theta(u) = \sum_{n \in \Z} q^{n(n-1)/2} u^n\) used in the construction of the Tate uniformization.
\end{description}
